% from ega4-9-fr.tex
\oldpage[IV-3]{71}

que l'anneau quotient $B/(T-\alpha)B$ admet par suite des éléments nilpotents~;
autrement dit, $X_s$ n'est réduit pour aucun $s\neq\eta$, tandis que
$X_\eta=\operatorname{Spec}(L)$ est intègre.

Nous verrons un peu plus loin (9.7) que l'on obtient par contre des propriétés
constructibles lorsqu'on considère les notions «~géométriques~»
correspondantes aux notions de préschéma irréductible, réduit ou intègre (4.5
et 4.6).
\end{remark}

\subsection{Constructibilité de certaines propriétés des morphismes}
\label{subsection:IV.9.6-fr}

\begin{proposition}[9.6.1]
\label{IV.9.6.1-fr}
Soient $X$, $Y$ deux $S$-préschémas de présentation finie, $f:X\to Y$ un
$S$-morphisme. Soit $E$ l'ensemble des $s\in S$ pour lesquels $f_s$ a l'une
des propriétés suivantes~: être~:
\begin{enumerate}[label=(\roman*),widest=xii]
  \item surjectif~;
  \item dominant~;
  \item séparé~;
  \item propre~;
  \item radiciel~;
  \item fini~;
  \item quasi-fini~;
  \item une immersion~;
  \item une immersion fermée~;
  \item une immersion ouverte~;
  \item un isomorphisme~;
  \item un monomorphisme.
\end{enumerate}
Alors $E$ est localement constructible dans $S$.
\end{proposition}

Les assertions de (i), (v) et (vii) ne sont insérées que pour mémoire, ayant
déjà été établies dans (9.3.2).

(ii)~: Notons d'abord que $f$ est de présentation finie (1.6.2, (v)), donc, en
vertu du théorème de Chevalley (1.8.4), $f(X)$ est localement constructible dans
$Y$. Par ailleurs, on a $f_s(X_s)=(f(X))_s$ (\textbf{I}, 3.6.1)~; l'ensemble
des $s$ tels que $f_s$ soit dominant est l'ensemble des $s$ tels que
$(f(X))_s$ soit dense dans $Y_s$~; la conclusion résulte donc dans ce cas de
(9.5.3).

(iii)~: Comme $f:X\to Y$ est de présentation finie, l'immersion diagonale
$\Delta_f:X\to X\times_YX$ est de présentation finie (1.6.2, (iv) et (v))~; il
résulte donc de (1.8.4.1) que $\Delta_f(X)=\Delta_X$ est une partie localement
constructible de $X\times_YX$. Dire que $f_s$ est séparé signifie (compte tenu
de (\textbf{I}, 5.3.4)) que $(\Delta_X)_s$ est fermé dans
$X_s\times_{Y_s}X_s=(X\times_YX)_s$~; la conclusion résulte donc ici de
(9.5.4).

(iv)~: Montrons que la propriété «~$X$ et $Y$ sont des préschémas algébriques
sur un corps $k$, et $f:X\to Y$ un $k$-morphisme propre~» est constructible. La
condition (9.2.1, (i)) est vérifiée en vertu de (2.7.1, (vii)). On peut donc se
borner au cas où $S$ est affine, noethérien et intègre, de point générique
$\eta$, et on doit prouver que $E$ ou $S-E$ est un voisinage de $\eta$.
Supposons d'abord que $\eta\in E$~; il résulte de (8.10.5, (xii)) appliqué
suivant la méthode de (8.1.2, \emph{a)}) qu'en remplaçant $S$ par un voisinage
de $\eta$, on peut

